% ----- CHAUPAI 12 -----

\newpage

\subsection*{\deva{द्वादश चौपाई} (Twelfth Chaupai)}
\addcontentsline{toc}{subsection}{\deva{द्वादश चौपाई} (Twelfth Chaupai) — \deva{रघुपति कीन्ही...}}

\begin{minipage}[t]{0.55\textwidth}
\vspace{0pt}

{\large \linespread{1.5}\selectfont
\deva{रघुपति कीन्ही बहुत बड़ाई।}\\
\deva{तुम मम प्रिय भरतहि सम भाई॥}
\par}

\vspace{0.5em}

{\large \linespread{1.5}\selectfont
\telu{రఘుపతి కీన్హీ బహుత బడాయీ।}\\
\telu{తుమ మమ ప్రియ భరతహి సమ భాయీ॥}
\par}

\vspace{0.5em}

{\large \linespread{1.3}\selectfont
\textit{raghupati kīnhī bahuta baṛāī |}\\
\textit{tuma mama priya bharatahi sama bhāī ||}
\par}
\end{minipage}%
\hfill
\begin{minipage}[t]{0.42\textwidth}
\vspace{0pt}
\begin{tcolorbox}[anchorbox]
\textbf{Recognition of Merit:} Lord Rama praised Hanuman by stating he was equal to his own royal brother Bharata. A true leader creates a meritocracy. They judge people by their character, loyalty, and hard work, not by their family background, species, or social status.
\vspace{0.5em}\newline The supreme leader publicly embraced and rewarded him, proving that merit and hard work transcend all social titles.\newline\null\hfill\href{https://www.valmiki.iitk.ac.in/sloka?field_kanda_tid=6&language=dv&field_sarga_value=1}{\textit{Yuddha Kanda, 1:12-14}}
\end{tcolorbox}
\end{minipage}

\vspace{0.5em}
\subsubsection*{\textbf{Visual Rhythm Map:}}
\begin{table}[h!]
\centering
\renewcommand{\arraystretch}{1.2}
\setlength{\tabcolsep}{0pt}
\begin{tabularx}{\textwidth}{|*{16}{>{\centering\arraybackslash\hsize=1.0\hsize}X|}}
\hline
\multicolumn{4}{|>{\centering\arraybackslash\hsize=4.0\hsize}X|}{\textbf{\laghu{र}\laghu{घु}\laghu{प}\laghu{ति}} \par (4)} &
\multicolumn{4}{>{\centering\arraybackslash\hsize=4.0\hsize}X|}{\textbf{\guru{कीं}\guru{न्ही}} \par (4)} &
\multicolumn{3}{>{\centering\arraybackslash\hsize=3.0\hsize}X|}{\textbf{\laghu{ब}\laghu{हु}\laghu{त}} \par (3)} &
\multicolumn{5}{>{\centering\arraybackslash\hsize=5.0\hsize}X|}{\textbf{\laghu{ब}\guru{ड़ा}\guru{ई}} \par (5)} \\
\hline
\end{tabularx}
\vspace{-1.5em}
\end{table}

\begin{table}[h!]
\centering
\renewcommand{\arraystretch}{1.2}
\setlength{\tabcolsep}{0pt}
\begin{tabularx}{\textwidth}{|*{16}{>{\centering\arraybackslash\hsize=1.0\hsize}X|}}
\hline
\multicolumn{2}{|>{\centering\arraybackslash\hsize=2.0\hsize}X|}{\textbf{\laghu{तु}\laghu{म}} \par (2)} &
\multicolumn{2}{>{\centering\arraybackslash\hsize=2.0\hsize}X|}{\textbf{\laghu{म}\laghu{म}} \par (2)} &
\multicolumn{2}{>{\centering\arraybackslash\hsize=2.0\hsize}X|}{\textbf{\laghu{प्रि}\laghu{य}} \par (2)} &
\multicolumn{4}{>{\centering\arraybackslash\hsize=4.0\hsize}X|}{\textbf{\laghu{भ}\laghu{र}\laghu{त}\laghu{हि}} \par (4)} &
\multicolumn{2}{>{\centering\arraybackslash\hsize=2.0\hsize}X|}{\textbf{\laghu{स}\laghu{म}} \par (2)} &
\multicolumn{4}{>{\centering\arraybackslash\hsize=4.0\hsize}X|}{\textbf{\guru{भा}\guru{ई}} \par (4)} \\
\hline
\end{tabularx}
\vspace{-1.5em}
\end{table}

\vspace{1em}
\textbf{ \deva{सरल-संस्कृतम्}:}\\
 \deva{श्रीरामः (रघुपतिः) भवतः अतीव प्रशंसां (बड़ाई) कृतवान्।}\\
 \deva{"त्वं मम (भरतेन समानः) प्रिय-भ्राता असि"॥}\\

The Lord of the Raghus praised you immensely.
He said, "You are as dear a brother to me as Bharata."


\vspace{1em}
\newpage
\textbf{\deva{प्रतिपदार्थः} (Meaning of each word):}
\begin{table}[h!]
\centering
\renewcommand{\arraystretch}{1.2}
\begin{tabularx}{\textwidth}{@{} l l X X @{}}
\toprule
\textbf{Awadhi} & \textbf{Sanskrit} & \textbf{English} & \textbf{Grammar \& Notes} \\
\midrule
\deva{रघुपति} / \telu{రఘుపతి}  &  \deva{रघुपतिः}  &  Lord Rama  &   \\
\deva{कीन्ही} / \telu{కీన్హీ}  &  \deva{कृतवान्}  &  Did  & Awadhi past tense verb \\
\deva{बहुत} / \telu{బహుత}  &  \deva{अतीव / बहु}  &  Much / Lots of  &   \\
\deva{बड़ाई} / \telu{బడాయీ}  &  \deva{प्रशंसा / महिमा}  &  Praise / Greatness  &   \\
\deva{तुम} / \telu{తుమ}  &  \deva{त्वम्}  &  You  &   \\
\deva{मम} / \telu{మమ}  &  \deva{मम}  &  My  & Pure Sanskrit! \\
\deva{प्रिय} / \telu{ప్రియ}  &  \deva{प्रियः}  &  Dear  & Pure Sanskrit! \\
\deva{भरतहि} / \telu{భరతహి}  &  \deva{भरताय / भरतेन}  &  Bharata  & Objective suffix `-hi` \\
\deva{सम} / \telu{సమ}  &  \deva{समानः}  &  Equal / Like  &   \\
\deva{भाई} / \telu{భాయీ}  &  \deva{भ्राता}  &  Brother  &   \\
\bottomrule
\end{tabularx}
\end{table}



\newpage
